\subsection{索伯列夫空间}

设 U 是 \(\mathbf{R}^{n} \) 的开子集。设 p 是实数 \(\geqslant 1\) ， k 是自然数。我们用 \(W^{k,p}(U)\) 表示满足如下条件的广义函数 \(f \in \mathcal{D}'(U)\) 所成的空间：对每个 \(\alpha \in N^{n}\) （满足 \(|\alpha| \leqslant k\) ），广义函数 \(\partial^{\alpha}f\) 与 \(\mathrm{L}^{p}(U) \) 的一个元素相伴。

特别地，当 \(U = \mathbf{R}^{n}\) 时， \(W^{k,p}(U)\) 中的元素是缓增广义函数。

从 \(\mathrm{W}^{k,p}(\mathrm{U})\) 到 \(\mathbf{R}_{+}\) 中的、把 \(f\in \mathrm{W}^{k,p}(\mathrm{U})\)

\[
\| f \| _ {k, p} = \Bigl (\sum_ {| \alpha | \leqslant k} \| \partial^ {\alpha} f \| _ {p} ^ {p} \Bigr) ^ {1 / p}
\]

是 \(W^{k,p}(U)\) 上的一个范数。空间 \(W^{k,p}(U)\) 将始终赋有这一范数；这一赋范空间称为指标为 k 、指数为 p 的索伯列夫空间。

空间 \(\mathcal{D}(\mathrm{U})\) 包含在 \(\mathrm{W}^{k,p}(\mathrm{U})\) 中。我们用 \(\mathrm{W}_0^{k,p}(\mathrm{U})\) 表示 \(\mathcal{D}(\mathrm{U})\) 在 \(\mathrm{W}^{k,p}(\mathrm{U})\) 中的闭包。它是 \(\mathrm{W}^{k,p}(\mathrm{U})\) 的一个闭子空间。

有 \(\mathrm{W}_{0}^{k,p}(\mathbf{R}^{n})=\mathrm{W}^{k,p}(\mathbf{R}^{n})\) ，但空间 \(\mathrm{W}^{k,p}(\mathrm{U})\) 与 \(\mathrm{W}_{0}^{k,p}(\mathrm{U})\) 一般是不同的（参见 IV, p. 334 的习题 12 与 IV, p. 334 的习题 14）。

我们也记 \(\mathrm{H}^k (\mathrm{U}) = \mathrm{W}^{k,2}(\mathrm{U})\) ， \(\mathrm{H}_0^k (\mathrm{U}) = \mathrm{W}_0^{k,2}(\mathrm{U})\) 。

 \(\mathrm{H}^{k}(\mathrm{U})\) 的范数是准希尔伯特范数，它由 \(\mathrm{H}^{k}(\mathrm{U})\) 上的、由下式定义的正埃尔米特形式所导出。

\[
\left(f _ {1}, f _ {2}\right) \mapsto \sum_ {| \alpha | \leqslant k} \int_ {\mathrm{U}} \overline {{\partial^ {\alpha} f _ {1}}} \partial^ {\alpha} f _ {2} d \mu .
\]

希尔伯特空间 \(\mathrm{H}^k (\mathrm{U})\) 与 EVT, V, p. 6 例 (3) 中记作 \(\mathcal{H}^k\) 的空间重合。

按定义有 \(\mathrm{W}^{0,p}(\mathrm{U})=\mathrm{L}^{p}(\mathrm{U})\) ， \(\mathrm{H}^{0}(\mathrm{U})=\mathrm{L}^{2}(\mathrm{U})\) ；此外，由 IV, p. 202 的命题 4, b)， \(\mathrm{W}_{0}^{0,p}(\mathrm{U})=\mathrm{L}^{p}(\mathrm{U})\) 。

命题 20. — 索伯列夫空间 \(W^{k,p}(U)\) 与 \(W_{0}^{k,p}(U)\) 是可数型的巴拿赫空间。特别地，空间 \(H^{k}(U)\) 与 \(H_{0}^{k}(U)\) 是可数型的希尔伯特空间。

只须证明关于 \(\mathrm{W}^{k,p}(\mathrm{U})\) 的断言。

设 I 是满足 \(\alpha \in \mathbf{N}^n\) 的 \(|\alpha| \leqslant k\) 所成的集。线性映射 \(u\) （从 \(\mathrm{W}^{k,p}(\mathrm{U})\) 到 \(\mathrm{L}^p (\mathrm{U})^{\mathrm{I}}\) 中的）把 \(f\) 对应到族 \((\partial^{\alpha}f)_{\alpha \in \mathrm{I}}\) ，它是单射；按 \(\mathrm{W}^{k,p}(\mathrm{U})\) 上范数的定义，它是连续且严格的。为证 \(\mathrm{W}^{k,p}(\mathrm{U})\) 完备，只须证它在 \(u\) 下的像是闭的。设 \((f_n)_{n \in \mathbf{N}}\) 是 \(\mathrm{W}^{k,p}(\mathrm{U})\) 中的序列，使得 \((u(f_n))_{n \in \mathbf{N}}\) 收敛。设 \((g_{\alpha})_{\alpha \in \mathrm{I}} \in \mathrm{L}^p (\mathrm{U})^{\mathrm{I}}\) 是它的极限。对 \(\alpha \in \mathrm{I}\) ，序列 \((\partial^{\alpha}f_n)_{n \in \mathbf{N}}\) 在 \(\mathrm{L}^p (\mathrm{U})\) 中收敛于 \(g_{\alpha}\) 。当然，收敛也在 \(\mathcal{D}'(\mathrm{U})\) 中发生。令 \(f = g_0\) 。对每个 \(\varphi \in \mathcal{D}(\mathrm{U})\) ，得到

\[
\begin{array}{r l} & {\langle \partial^ {\alpha} f, \varphi \rangle = (- 1) ^ {| \alpha |} \langle f, \partial^ {\alpha} \varphi \rangle} \\ & {\qquad = \lim _ {n \to + \infty} (- 1) ^ {| \alpha |} \langle f _ {n}, \partial^ {\alpha} \varphi \rangle = \lim _ {n \to + \infty} \langle \partial^ {\alpha} f _ {n}, \varphi \rangle = \langle g _ {\alpha}, \varphi \rangle .} \end{array}
\]

这就证明了 \(g_{\alpha} = \partial^{\alpha} f\) 对所有 \(\alpha \in I\) 成立，从而 \((g_{\alpha})_{\alpha \in I} = u(f)\) 属于 u 的像。

空间 \(\mathrm{W}^{k,p}(\mathrm{U})\) 由 u 等同于空间 \(\mathrm{L}^{p}(\mathrm{U})^{\mathrm{I}}\) 的一个子空间；后者是可数型的（IV, p. 180 的命题 2 与 TG, IX, p. 19，cor., (ii))），因而 \(\mathrm{W}^{k,p}(\mathrm{U})\) 也是可数型的（loc. cit., (i)）。

命题 21. --- 设 N 是 \(\mathbf{R}^{n}\) 上的欧几里得范数。设 \(k\) 是整数 \(\geqslant 0\) 。索伯列夫空间 \(\mathrm{H}^{k}(\mathbf{R}^{n})\) 是这样的 \(f \in \mathcal{S}'(\mathbf{R}^{n})\) 所成的空间： \((1 + \mathrm{N}^{k})\mathcal{F}(f)\) 属于 \(\mathrm{L}^{2}(\mathbf{R}^{n})\) 。

我们对 \(k\) 作归纳。当 \(k = 0\) 时，结论由普朗谢雷尔定理（II, p. 215，定理 1）得出。设 \(k = 1\) 。对 \(f \in \mathcal{S}'(\mathbf{R}^n)\) ，我们有 \((1 + \mathrm{N})\mathcal{F}f \in \mathrm{L}^2(\mathbf{R}^n)\) 当且仅当 \(f \in \mathrm{L}^2(\mathbf{R}^n)\) 且 \(\mathrm{N}\mathcal{F}f \in \mathrm{L}^2(\mathbf{R}^n)\) 。此外， \(\mathrm{N}\mathcal{F} \in \mathrm{L}^2(\mathbf{R}^n)\) 当且仅当：对每个 \(\alpha \in \mathbf{N}^n\) （满足 \(|\alpha| = 1\) ），有 \(\mathrm{X}^\alpha\mathcal{F}f \in \mathrm{L}^2(\mathbf{R}^n)\) 。由于 \(\mathcal{F}(\partial^\alpha f) = 2i\pi\mathrm{X}^\alpha\mathcal{F}f\) ，这一条件意味着对每个 \(\mathcal{F}(\partial^\alpha f) \in \mathrm{L}^2(\mathbf{R}^n)\) （满足 \(\alpha\) ）有 \(|\alpha| = 1\) ，亦即对每个 \(\partial^\alpha f \in \mathrm{L}^2(\mathbf{R}^n)\) （满足 \(\alpha\) ）有 \(|\alpha| = 1\) 。由此得出断言对 \(k = 1\) 为真。

现在设 \(k \geqslant 2\) ，并且关于 \(\mathrm{H}^{\ell}(\mathbf{R}^{n})\) 的断言对每个正整数 \(\ell \leqslant k - 1\) 成立。设 \(f \in \mathcal{S}'(\mathbf{R}^{n})\) 。按定义， \(f \in \mathrm{H}^{k}(\mathbf{R}^{n})\) 当且仅当 \(f \in \mathrm{L}^{2}(\mathbf{R}^{n})\) ，且对每个 \(\beta \in \mathbf{N}^{n}\) （满足 \(|\beta| \leqslant 1\) ），广义函数 \(\partial^{\beta} f\) 属于 \(\mathrm{H}^{k-1}(\mathbf{R}^{n})\) 。由归纳假设，这等价于 \(f \in \mathrm{L}^{2}(\mathbf{R}^{n})\) ，且对每个 \((1 + \mathrm{N}^{k-1})\mathcal{F}(\partial^{\beta} f) \in \mathrm{L}^{2}(\mathbf{R}^{n})\) （满足 \(\beta \in \mathbf{N}^{n}\) ）有 \(|\beta| \leqslant 1\) 。由于 \(\mathcal{F}(\partial^{\beta} f) = (2i\pi X)^{|\beta|}\mathcal{F}(f)\) ，条件 \(f \in \mathrm{H}^{k}(\mathbf{R}^{n})\) 等价于说： \(\mathcal{F}f \in \mathrm{L}^{2}(\mathbf{R}^{n})\) ，且对每个 \((1 + \mathrm{N}^{k-1})\mathrm{X}^{\beta}\mathcal{F}f \in \mathrm{L}^{2}(\mathbf{R}^{n})\) （满足 \(\beta \in \mathbf{N}^{n}\) ）有 \(|\beta| \leqslant 1\) 。

不等式

\[
1 + \mathrm{N}^{k}\leqslant 1 + \mathrm{N}^{k - 1}\sum_{\substack{\beta \in \mathbf{N}^{n}\\ |\beta |\leqslant 1}}|\mathrm{X}^{\beta}|\leqslant 1 + n^{1 / 2}\mathrm{N}^{k}\leqslant (1 + n^{1 / 2})(1 + \mathrm{N}^{k})
\]

于是蕴含 \(f \in \mathrm{H}^k(\mathbf{R}^n)\) 当且仅当 \((1 + \mathrm{N}^k)\mathcal{F} \in \mathrm{L}^2(\mathbf{R}^n)\) 。

